\documentclass[12pt, a4paper]{article}

\usepackage[T1]{fontenc}
\usepackage[utf8]{inputenc}
\usepackage{palatino}
\usepackage{microtype}
\usepackage[margin=1.25in]{geometry}
\usepackage{setspace}
\usepackage{parskip}
\usepackage{titlesec}
\usepackage[numbers]{natbib}
\usepackage{hyperref}
\usepackage{amsmath}

\hypersetup{
  colorlinks=true,
  linkcolor=black,
  citecolor=black,
  urlcolor=black
}

\titleformat{\section}{\normalfont\large\scshape}{\thesection.}{1em}{}
\titleformat{\subsection}{\normalfont\normalsize\itshape}{\thesubsection.}{1em}{}

\setstretch{1.15}

\title{\textsc{Civilization's Undo Stack}\\
\large On Markov Boundaries, Coordination Costs, and\\
the Recurring Regression of Progress}
\author{Flyxion}
\date{June 2026}

\begin{document}

\maketitle

\begin{abstract}
Across apparently unrelated domains---interface design, computational
architecture, archival systems, infrastructure, governance, and
artificial intelligence---a common structural failure recurs with
remarkable consistency: mechanisms developed specifically to absorb
coordination costs are dismantled, and the costs are re-externalized
onto individuals. This essay argues that the underlying concept
unifying these failures is the Markov boundary: a structure that
creates conditional independence between phases of activity, preventing
errors from propagating indefinitely. Trash bins, cooling-off periods,
appeals courts, peer review cycles, undo stacks, and NULL wavefronts
are all, at the appropriate level of abstraction, Markov boundaries.
Their removal is always presented as simplification. Its actual
effect is the indefinite propagation of errors across episodes that
were previously separable. A mature civilization is one that builds
Markov boundaries. The recurring temptation to remove them in the
name of efficiency, clarity, or novelty constitutes a specific and
diagnosable failure mode of institutional intelligence.
\end{abstract}

\section{The Problem of Propagating Error}

Not all mistakes stay where they land.

Some errors are local. A misplaced word is corrected by the next
sentence. A wrong turn is corrected by the next intersection. A
miscalculation is caught by the next check. The error occurs, it is
noticed, it is repaired, and the work continues from approximately
where it would have been had the error not occurred at all. The cost
is bounded. It does not accumulate.

Other errors propagate. A contaminated water supply does not merely
harm the person who drinks first; it flows downstream. A mistaken
assumption early in a long calculation does not merely corrupt one
step; it invalidates every step that follows. A decision made in panic
does not merely affect the moment of decision; it forecloses options
for months or years afterward. The error enters a system and travels
through it, touching everything it passes.

The difference between these two kinds of error is not in the magnitude
of the initial mistake. It is in the structure of the system through
which the mistake moves. Some systems are built to contain errors.
Others are built in ways that allow errors to travel freely---sometimes
by oversight, sometimes by design, sometimes by the gradual removal
of structures that once contained them.

The history of civilization is substantially the history of learning,
through painful experience, which errors propagate and why---and then
building structures to stop them.

\section{What a Markov Boundary Is}

In probability theory, a Markov property holds when the future of a
system is conditionally independent of its past, given its present
state. Knowing the present state, additional knowledge of the history
provides no further information about what comes next. The present is
a sufficient summary. History can be safely ignored beyond the
boundary.

This is an abstract mathematical statement, but it describes something
concrete and important: a boundary that separates episodes from each
other in a way that prevents what happened in one episode from
contaminating the next. The boundary does not eliminate history. It
determines what must cross from one episode to the next, and what can
be safely discarded.

A Markov boundary, as this essay uses the term, is any structure that
achieves this separation in practice. It need not be mathematically
exact. It need not be formally specified. It is any mechanism that
prevents errors from propagating indefinitely from one phase of
activity into the next by creating a clean interface between episodes,
determining what information must cross and what can be absorbed.

The key property is conditional independence: given the state of the
boundary, what comes after is independent of how things went before.
The error may have been catastrophic. But if it was absorbed at the
boundary, the next episode begins clean.

This gives us a four-level structure that the rest of this essay
develops. Markov boundaries create conditional independence between
episodes. Conditional independence absorbs coordination costs---the
costs that would otherwise be borne by whoever must manage the
transition between phases. Civilizations progress, in part, by
building structures that absorb these costs into infrastructure rather
than externalising them onto individuals. And civilizations regress,
in a specific and diagnosable way, when those structures are removed
and the costs are pushed back out. The essays that follow this one
each examine one domain where that regression has occurred.

\section{The Invisible Infrastructure of Separation}

Once the concept is named, Markov boundaries appear everywhere---which
is exactly what one should expect if civilizations that survived
learned to build them.

A trash bin is a temporal Markov boundary. It separates the episode
of \emph{deciding to delete} from the episode of \emph{deletion
becoming permanent}. Given the current state of the trash bin, the
next episode---recovery or confirmation---is independent of the
circumstances under which the deletion was initiated. Whether the
deletion was deliberate, accidental, fatigued, or panicked does not
matter. The boundary absorbs all of those into a uniform interface:
the content is here, and it is recoverable.

A cooling-off period is a temporal Markov boundary of a different
kind. A decision made in anger is not automatically final. The period
of enforced delay separates the episode of emotional arousal from the
episode of committed action. Given the state at the end of the
cooling-off period---the person still wishes to proceed---the action
follows. The emotional history of how the decision was reached does
not propagate. It has been absorbed.

An appeals court is a legal Markov boundary. It separates the episode
of a trial, with all its contingencies of evidence, argument, fatigue,
and bias, from the episode of irreversible sentence. The conviction
may stand or it may not; what matters is that the process of
conviction is separated from the process of finalization by a
boundary that absorbs the errors of the first before permitting entry
to the second.

Peer review is an epistemic Markov boundary. It separates the episode
of individual inquiry, which is prone to confirmation bias, motivated
reasoning, and simple error, from the episode of knowledge
entering the shared record. The content of a paper's history---how
many times it was revised, how much the author struggled, whether the
experiments were conducted on a bad day---does not propagate into the
status of the published result. The boundary absorbs all of that. What
crosses is the output of peer evaluation, not the history of
production.

Version control is a computational Markov boundary. A Git commit
separates the episode of development from the episode of recorded
history. The chaos of how the code was written---the false starts, the
deleted experiments, the working-at-3am decisions---does not propagate
into the commit. What crosses is the state of the codebase at the
moment of deliberate commitment. Every subsequent operation on the
repository is independent of the development history except insofar
as that history is explicitly encoded in the commits themselves.

Sleep is a neurological Markov boundary. The consolidation processes
of sleep separate the episode of experience from the episode of
memory. Information that entered the brain during the day is processed,
evaluated, and either integrated or discarded. The dream state is not
accessible to waking cognition, which means waking cognition begins
each morning with a compressed and partially cleaned record of the
previous day rather than the full unfiltered stream of experience.
Chronic sleep deprivation produces, among other things, exactly what
one would expect from degraded Markov boundaries: errors that would
ordinarily be resolved during sleep persist and accumulate, affecting
subsequent reasoning in ways that are difficult to trace.

In each case the structure is the same. A phase of activity generates
errors, contingencies, and noise. The Markov boundary absorbs those
into a defined interface. The next phase begins from the boundary
state, not from the full history of everything that preceded it. The
cost of coordination between phases is absorbed by the boundary itself
rather than externalised onto whoever must manage the transition.

\section{The Removal Pattern}

Markov boundaries are expensive to build and maintain. They require
infrastructure: engineers to implement the trash bin, institutions
to staff the appeals court, reviewers to conduct peer review, time to
allow for the cooling-off period. They introduce delay. They consume
resources. They add complexity to systems that might otherwise appear
simpler without them.

This makes them permanently vulnerable to a specific argument: that
they are unnecessary, that they represent over-engineering or
bureaucratic excess, that removing them would improve efficiency, that
whatever they protect against is not a real problem, or that a newer
paradigm has rendered them obsolete.

The argument is almost always wrong. And it is almost always made in
good faith.

The pattern of removal follows a recognisable sequence. A boundary
is built to solve a specific, painful problem. Over time, the boundary
works: the problem ceases to be painful, because the boundary is
absorbing it. A new generation of practitioners encounters the
boundary without having encountered the original problem. To them, the
boundary appears as overhead---a cost without visible benefit. They
propose its removal in the name of simplicity. The removal is
implemented. The original problem reappears. The cycle begins again.

This is not stupidity. It is a structural consequence of successful
engineering. A boundary that works perfectly is invisible. Invisible
things do not generate advocates. They are removed by people who never
saw the problem they solved.

The additional compounding factor is that the benefits of a Markov
boundary are often diffuse and statistical while its costs are
concentrated and immediate. The trash bin costs storage and interface
complexity today. Its benefit is recovery from an accidental deletion
that may not occur for months---and when it does occur, the person
affected may not be the person who made the removal decision. The
costs and benefits are borne by different people at different times,
which makes the trade-off systematically illegible to whoever is
optimising the system at any given moment.

\section{Removing Boundaries Is Always Called Something Else}

The removal of Markov boundaries is never described as the removal of
Markov boundaries. It is always described in terms that make the
removal appear as an improvement.

Removing the trash bin is called \emph{clarity}: deletion means
deletion; there is no ambiguity; the system is honest about what
irreversibility means. The moral language of transparency is
recruited to justify the elimination of a safety mechanism.

Removing cooling-off periods is called \emph{responsiveness}: the
system reacts immediately to user intent; there is no artificial
delay; the experience is frictionless. The language of service quality
is recruited to justify the elimination of temporal buffering.

Removing appeals processes is called \emph{efficiency}: cases are
resolved faster; the system is not clogged with endless re-litigation;
finality is achieved promptly. The language of throughput is recruited
to justify the elimination of error-correction.

Removing peer review---or replacing it with citation counts, download
statistics, and social media engagement---is called
\emph{democratisation}: knowledge is no longer gatekept by insular
elites; the public can evaluate claims directly; science is made
accessible. The language of equality is recruited to justify the
elimination of epistemic filtration.

In every case the rhetorical move is the same: the Markov boundary
is recharacterised as an obstacle rather than infrastructure. The
costs of the boundary are made visible; its function is made
invisible. This recharacterisation is not always dishonest.
Sometimes it reflects genuine ignorance of what the boundary does.
More often it reflects an optimization pressure that is indifferent
to the boundary's function---optimising for speed, engagement,
revenue, or apparent simplicity, and treating error absorption as
someone else's concern.

\section{Re-Externalisation and Its Costs}

When a Markov boundary is removed, the coordination cost it was
absorbing does not disappear. It is re-externalised. Someone still
pays it. The question is only who.

When the trash bin is removed, the user who accidentally deletes
months of work pays the full cost. When the cooling-off period is
removed, the person who sends an irretrievable message in a moment of
anger pays the full cost. When the appeals process is shortened, the
person wrongly convicted pays the full cost. When peer review is
replaced by metrics, the reader who cannot distinguish rigorous
research from motivated confabulation pays the full cost.

The system appears more efficient because it has externalised the cost
onto individuals. The cost has not been reduced. It has been
redistributed from the system to the user, from the organisation to
the person with the least power to absorb it.

This redistribution tends to be regressive. Markov boundaries, when
they function well, absorb errors regardless of the resources of the
person who made them. A trash bin recovers a file whether its owner
is a software engineer who would have noticed the loss immediately or
a first-time user who might not notice until the work was urgently
needed. An appeals court provides recourse regardless of whether the
defendant has the social capital to navigate the consequences of a
wrongful conviction. Peer review filters out low-quality work
regardless of the institutional prestige of its author.

When those boundaries are removed and costs are re-externalised, the
resulting distribution of harm follows the prior distribution of
resources. Those with the knowledge, time, and technical capacity to
build personal compensating mechanisms---manual backups, private
copies, duplicate records---absorb the cost with relatively little
damage. Those without those resources bear the full consequence of
errors that the boundary would have absorbed.

Re-externalisation is not merely inefficient. It is structurally
inequitable.

\section{The Specific Intelligence of Boundary-Building}

Recognising that Markov boundaries are valuable does not itself
explain why they are difficult to build or why their loss is
difficult to reverse. The difficulty is not technical. It is
epistemological.

Building a Markov boundary requires knowing, in advance, which errors
will occur and at what rate, what information must cross the boundary
and what can be absorbed, how much delay the separation requires, and
what the cost of false positives---boundaries that block valid
transitions---will be. This knowledge is not available a priori. It
is accumulated through experience of failure: through seeing what
happens when errors propagate, through measuring the distribution of
mistake types, through learning which transitions genuinely need
buffering and which do not.

Boundary-building is therefore an empirical activity, and the
empirical record from which it draws is the record of past failures.
A society that has experienced catastrophic water contamination knows
to separate waste flows from drinking water. A discipline that has
experienced irreproducible results knows to require replication. A
legal system that has experienced systematic wrongful conviction knows
to provide appeals. The boundary encodes accumulated knowledge about
where errors arise, how they propagate, and what it costs to absorb
them rather than let them travel.

This is what makes the removal of boundaries epistemically dangerous
in a way that goes beyond the immediate cost of any single failure.
A boundary that is removed takes with it the knowledge encoded in its
construction. The institutional memory of why the boundary existed is
not automatically preserved when the boundary disappears. New
practitioners face the same distribution of errors as their
predecessors but without the infrastructure those predecessors built
to absorb them---and, often, without even knowing that such
infrastructure once existed.

Civilisational knowledge can be destroyed not only by catastrophe but
by the quiet removal of the infrastructure that embodied it.

\section{The General Thesis}

The essays that follow this one examine a range of domains in which
the same pattern appears: a coordination problem was identified; a
Markov boundary was developed to absorb its costs; the boundary was
removed or allowed to degrade; the original problem reappeared,
usually under a different name.

In interface design, the boundary is the staged deletion
architecture---the undo stack, the trash bin, the temporal latency
buffer between intention and irreversibility. Its removal is justified
by appeals to simplicity and the moral language of honest finality.
But its true cost is the collapse of the conditional independence
between an ambiguous human intention and a permanent state change.
The coordination cost of perfectly synchronising current behaviour
with future necessity is re-externalised onto the user, whose
admissible future trajectories are truncated by a single uninsulated
keystroke.

In computational architecture, the boundary is the propagating
completion signal---the NULL wavefront, the intermediate
representation, the process chain's firebreak between episodes. Its
removal is the logical consequence of an application-centric paradigm
that localises state management in the user rather than in the
execution structure. Its cost is the dissolution of the separations
that prevent error cascades across asynchronous processes. Without a
completion signal to re-establish conditional independence between
distinct computational episodes, the user becomes the manual
coordinator of every intermediate state the system once absorbed
automatically.

In archival systems, the boundary is the finding aid, the metadata
schema, the retrieval pathway. Its removal---or rather, its
degradation under the weight of search-engine optimisation, link rot,
and algorithmic curation---is justified by the apparent sufficiency
of keyword search. Its cost is the loss of reachability: information
that is nominally preserved but practically irretrievable.

In infrastructure, the boundary is the local repair capacity, the
redundant system, the maintained buffer stock. Its removal is
justified by just-in-time efficiency. Its cost is the
re-externalisation of resilience onto communities that lack the
resources to build private substitutes for the public infrastructure
that was withdrawn.

In governance, the boundary is the legislative delay, the
deliberative procedure, the separation of powers. Its removal is
justified by the need for decisive action. Its cost is the
propagation of errors in judgment that deliberative procedures were
designed to catch.

In artificial intelligence, the boundary is the principled refusal,
the stable evaluative framework, the resistance to audience pressure.
Its removal is the consequence of optimising for engagement rather
than truth. Its cost is the propagation of whatever moral posture the
user arrives with, uncorrected, into their understanding of the world.

These are not analogies. They are instances of a single phenomenon,
and the phenomenon can be stated precisely: a Markov boundary is a
mechanism that prevents the state-space of one episode from becoming
the state-space of every subsequent episode. Without the boundary, a
local error does not stay local. It enters the global state and
remains there, contaminating every transition that follows. With the
boundary, the episode closes. The next episode begins from the
boundary state, not from the accumulated history of everything that
went wrong before it.

Each of the domains surveyed above involves the removal of exactly
this mechanism, and the re-externalisation of those costs onto
individuals.

\section{What Boundaries Reveal About Civilisation}

There is a temptation to interpret the recurring removal of Markov
boundaries as evidence of civilisational decline---of institutions
forgetting what they once knew. This interpretation is too simple.

The removal of boundaries is not a failure of intelligence. It is the
predictable outcome of a mismatch between the timescale on which
boundaries are built and the timescale on which organisations
optimise. Boundaries are built on the timescale of accumulated
experience, often spanning generations. Organisations optimise on
the timescale of quarterly performance, product releases, or
political cycles. From the perspective of the optimising organisation,
the Markov boundary looks like overhead. From the perspective of the
accumulated experience that built it, the boundary looks like
knowledge.

What the pattern reveals is not that civilisations are stupid. It is
that civilisations have a specific cognitive vulnerability: they lose
track of the difference between the cost of a boundary and the cost of
removing it. The cost of the boundary is immediate, visible, and
concentrated. The cost of its removal is deferred, statistical, and
distributed. These two kinds of cost are not commensurable in the
accounting systems that most organisations use. Optimising for the
first while ignoring the second is not a failure of moral reasoning.
It is a failure of accounting---specifically, the failure to
distinguish between the reduction of waste and the liquidation of
insurance.

A more robust civilisational intelligence would have two properties.
First, it would be capable of recognising when a cost is being
absorbed by infrastructure versus when it is being re-externalised
onto individuals. These look similar from the outside---in both cases,
the organisation appears to function without paying the cost---but
they are structurally opposite. Second, it would treat the removal of
a Markov boundary as evidence that requires justification proportional
to the difficulty of rebuilding the boundary, not proportional to the
apparent cost savings of the removal.

Neither property is currently standard. Both are achievable. The
first requires the conceptual vocabulary to name what is happening
when a boundary is removed. That is what this essay attempts to
provide.

\section{Building as a Form of Understanding}

There is a final dimension to the Markov boundary argument that goes
beyond the diagnostic.

A Markov boundary is not an arbitrary partition, nor a passive wall.
It is an active, crystallised hypothesis about the specific failure
modes of a system. To build a boundary that successfully introduces
conditional independence between two phases of activity, a designer
must explicitly confront the underlying error distribution of the
environment. One cannot design a trash bin without deciding what
constitutes a temporary state versus a permanent choice. One cannot
construct an electrical circuit breaker without calculating the
precise thermal threshold where current transforms from utility into
hazard. One cannot establish an appellate process without defining
what classes of trial-level deviation constitute systemic injustice
rather than tolerable variance.

The act of boundary-building therefore forces an adversarial
interrogation of the infrastructure that would otherwise never occur.
It requires the designer to answer three questions that the system
cannot answer about itself: What is the maximum volume of noise this
local episode can generate before it threatens the stability of the
whole? What is the minimum necessary information required to seed the
subsequent phase without dragging the historical baggage of the first
across the threshold? And where does the boundary itself fail? To
answer these questions rigorously is to generate knowledge. The
completed boundary becomes an operational archive of empirical
discoveries made during past collapses---a physical record of where
systems break, how errors propagate, and what it costs to contain
them.

This reveals the deepest cost of institutional amnesia. When an
optimising organisation removes a Markov boundary in the name of
efficiency, it does not merely re-externalise a coordination cost
onto the end-user. It destroys the epistemological record that made
the system legible to itself. New generations of practitioners are
left to navigate the exact same distribution of structural errors as
their predecessors, but they must do so stripped of both the
protection the boundary provided and the analytical vocabulary
required to understand why they are failing. Landauer showed that
erasing information generates thermodynamic heat---a physical cost
that cannot be avoided, only displaced \cite{landauer1961}. Erasing a
coordination boundary generates the cognitive equivalent: waste heat
that the individual must radiate, alone, in the form of vigilance,
defensive workarounds, and the repeated rediscovery of failure modes
that previous generations had already resolved and encoded.

The loss is not only of the infrastructure. It is of the knowledge
the infrastructure embodied.

The essays that follow examine what that loss looks like across
different domains. They are not primarily arguments for nostalgia.
They are attempts to make the boundary visible again---to name what
was removed, explain why it existed, and describe what its absence
costs.

That is the minimal first step. You cannot rebuild what you have
forgotten you lost.

\bigskip
\noindent\rule{\textwidth}{0.4pt}

\begin{thebibliography}{99}

\bibitem{markov1906}
A.~A. Markov.
\newblock Extension of the law of large numbers to dependent quantities.
\newblock \emph{Izvestiia Fiz.-Matem. Obshchestva pri Kazanskom Universitete},
  2nd ser., 15:135--156, 1906.

\bibitem{pearl1988}
J.~Pearl.
\newblock \emph{Probabilistic Reasoning in Intelligent Systems: Networks of
  Plausible Inference}.
\newblock Morgan Kaufmann, San Mateo, CA, 1988.

\bibitem{fant2005}
K.~M. Fant.
\newblock \emph{Logically Determined Design: Clockless System Design with NULL
  Convention Logic}.
\newblock John Wiley \& Sons, Hoboken, NJ, 2005.

\bibitem{fant2007}
K.~M. Fant.
\newblock \emph{Computer Science Reconsidered: The Invocation Model of Process
  Expression}.
\newblock John Wiley \& Sons, Hoboken, NJ, 2007.

\bibitem{norman2013}
D.~A. Norman.
\newblock \emph{The Design of Everyday Things}.
\newblock Basic Books, New York, 2013.

\bibitem{chesterton1929}
G.~K. Chesterton.
\newblock \emph{The Thing}.
\newblock Sheed \& Ward, London, 1929.

\bibitem{ostrom1990}
E.~Ostrom.
\newblock \emph{Governing the Commons: The Evolution of Institutions for
  Collective Action}.
\newblock Cambridge University Press, Cambridge, 1990.

\bibitem{scott1998}
J.~C. Scott.
\newblock \emph{Seeing Like a State: How Certain Schemes to Improve the Human
  Condition Have Failed}.
\newblock Yale University Press, New Haven, CT, 1998.

\bibitem{winner1980}
L.~Winner.
\newblock Do artifacts have politics?
\newblock \emph{Daedalus}, 109(1):121--136, 1980.

\bibitem{landauer1961}
R.~Landauer.
\newblock Irreversibility and heat generation in the computing process.
\newblock \emph{IBM Journal of Research and Development}, 5(3):183--191, 1961.

\bibitem{rawls1971}
J.~Rawls.
\newblock \emph{A Theory of Justice}.
\newblock Harvard University Press, Cambridge, MA, 1971.

\bibitem{illich1973}
I.~Illich.
\newblock \emph{Tools for Conviviality}.
\newblock Harper \& Row, New York, 1973.

\end{thebibliography}

\end{document}
